\documentclass[11pt]{article}
\usepackage[T1]{fontenc}
\usepackage{textcomp}
\usepackage[margin=0.75in]{geometry}
\usepackage{amsmath,amssymb,amsthm}
\usepackage{array,booktabs}
\usepackage{hyperref}
\setlength{\parindent}{0pt}
\newtheorem{theorem}{Theorem}
\theoremstyle{definition}
\newtheorem{definition}{Definition}
\title{Flow Proof Artifact: Comb-choose-succ-nineteen}
\author{Generated by \texttt{flow doc proof}}
\date{Flow Proof Book}
\begin{document}
\maketitle
Choosing nineteen from a successor set.\\[0.5em]
\textbf{Source.} graham-knuth-patashnik — *Concrete Mathematics*\\[0.5em]
\bigskip
\noindent{\large \textbf{Derived fact 1.} \textit{choose(succ(n), 19) = choose(n, 18) + choose(n, 19)} \hfill Comb \cdot choose \cdot choosing nineteen from a successor}\par
\smallskip\noindent\textit{Coordinate: Comb \cdot choose \cdot choosing nineteen from a successor}\par
\smallskip\noindent\textit{Source: graham-knuth-patashnik}\par
\smallskip\noindent\textit{Built on: pascal recurrence holds, for choose on Comb, choosing eighteen from a successor, for choose on Comb}\par
\medskip
\noindent\textbf{Goal.} choose(succ(n), 19) = choose(n, 18) + choose(n, 19)\par
\noindent$\displaystyle \forall n \in \mathbb{N}\quad choose(\mathrm{succ}(n), 19) = choose(n, 18) + choose(n, 19)$\par
\medskip
\noindent
\begin{tabular}{@{} >{\raggedright\arraybackslash}p{0.06\textwidth} >{\raggedright\arraybackslash}p{0.47\textwidth} >{\raggedleft\arraybackslash}p{0.41\textwidth} @{}}
\textbf{\#} & \textbf{Proof} & \textbf{Mathematics} \\
\midrule
\textbf{1.} & We prove that choosing nineteen from a successor for choose on Comb. & \\[0.45em]
\textbf{2.} & We invoke the definitional clause governing choose on Comb: pascal recurrence holds, for choose on Comb (instantiated for succ(n), 19). & \\[0.45em]
\textbf{3.} & We invoke the derived fact governing choose on Comb: choosing eighteen from a successor, for choose on Comb (instantiated for n). & \\[0.45em]
\textbf{4.} & From step 2 and step 3, this implies choose(the successor of n, 19) equals choose(n, 18) plus choose(n, 19). Hence proven. & $\displaystyle choose(\mathrm{succ}(n), 19) = choose(n, 18) + choose(n, 19)$ \\[0.45em]
\bottomrule
\end{tabular}
\label{thm:Comb-choose-choosing_nineteen_from_a_successor}
\par\medskip\hrule
\end{document}
